% !TeX program = XeLaTeX
% !TeX root = ../pujavidhanam.tex
\chapt{कृष्णाङ्गारक-चतुर्दशी-यम-तर्पणम्}

\twolineshloka*
{दीपोत्सवचतुर्दश्यां कार्यं तु यमतर्पणम्}
{कृष्णाङ्गारचतुर्दश्याम् अपि कार्यं सदैव वा}

\twolineshloka*
{कृष्णपक्षे चतुर्दश्याम् अङ्गारकदिनं यदा}
{तदा स्नात्वा शुभे तोये कुर्वीत यमतर्पणम्}

{\sffamily According to the above verses from Vaidyanātha Dīkṣitīyam (Āhnika Kāṇḍa, Uttarārdha), Yama Tarpaṇam must be performed on Naraka Chaturdashi (Dīpāvali), and also on Kṛṣṇāṅgāraka Chaturdashi, i.e. whenever the Chaturdashi of Kṛṣṇa Pakṣa falls on a Tuesday.}

\centerline{\bfseries जीवत्पिताऽपि कुर्वीत तर्पणं यमभीष्मयोः}\nobreak

{\sffamily As per the above vachana, Yama Tarpaṇam and Bhīṣma Tarpaṇam must be performed even by those whose father is alive.}

\twolineshloka*
{एकैकेन तिलैर्मिश्रान् दद्यात् त्रींस्त्रीन् जलाञ्जलीन्}
{संवत्सरकृतं पापं तत्क्षणादेव नश्यति}

\twolineshloka*
{कृष्णपक्षे चतुर्दश्यां यां काञ्चित् सरितं प्रति}
{यमुनायां विशेषेण नियतस्तर्पयेद् यमम्}

\twolineshloka*
{यत्र क्वचन नद्यां हि स्नात्वा कृष्णचतुर्दशीम्}
{सन्तर्प्य धर्मराजं तु मुच्यते सर्वकिल्बिषैः}

\twolineshloka*
{दक्षिणाभिमुखो भूत्वा तिलैः सव्यं समाहितः}
{देवतीर्थेन देवत्वात् तिलैः प्रेताधिपो यतः}

{\sffamily On such Kṛṣṇa Pakṣa Chaturdashi days, one who bathes in the Yamunā or any other river and offers Tarpaṇam to Dharmarāja (Yama) is instantly freed of all the pāpam accumulated over the entire year — such is the lofty phala described for Yama Tarpaṇam.}

\faClockO{} {\sffamily \textbf{Method:} Facing south, and taking black sesame (tila) along with the water, offer Tarpaṇam via Deva Tīrtha (the same tīrtha used for Deva Tarpaṇam in Sandhyāvandanam / Brahmayajñam), reciting each of the following names and offering water thrice for each.}

\centerline{\bfseries यज्ञोपवीतिना कार्यं प्राचीनावीतिनाऽथवा}

{\sffamily As per the above vachana, this Tarpaṇam may be performed wearing the Yajñopavītam either as Upavītī, or as Prācīnāvītī.}

\medskip

\dnsub{सङ्कल्पः}

\shuklambaradharam

\pranayama

ममोपात्त-समस्त-दुरितक्षयद्वारा श्रीपरमेश्वर-प्रीत्यर्थं शुभे शोभने मुहूर्ते अद्य ब्रह्मणः
द्वितीयपरार्धे श्वेतवराहकल्पे वैवस्वतमन्वन्तरे अष्टाविंशतितमे कलियुगे प्रथमे पादे
जम्बूद्वीपे भारतवर्षे भरत-खण्डे मेरोः दक्षिणे पार्श्वे शकाब्दे अस्मिन् वर्तमाने व्यावहारिकाणां
प्रभवादीनां षष्ट्याः संवत्सराणां मध्ये
\textbf{\blank\see{app:samvatsara_names}} नाम संवत्सरे
\textbf{\ayane}
\textbf{\rtu}-ऋतौ
\textbf{\masa}-मासे
\textbf{कृष्ण}-पक्षे
\textbf{चतुर्दश्यां} शुभतिथौ
\textbf{भौम}-वासरयुक्तायां
\textbf{\nakshatra}-नक्षत्र-%
\textbf{\yoga}-योग-%
\textbf{\karana}-करण-युक्तायां
च एवं गुण-विशेषण-विशिष्टायाम् अस्याम्\\
\textbf{चतुर्दश्यां} शुभतिथौ

ममोपात्त-समस्त-दुरितक्षयद्वारा श्रीपरमेश्वर-प्रीत्यर्थं यमधर्मराज-प्रीत्यर्थं कृष्णाङ्गारक-चतुर्दशी-पुण्यकाले यमतर्पणं करिष्ये।

\dnsub{तर्पण-मन्त्राः}

१. यमं तर्पयामि। यमं तर्पयामि। यमं तर्पयामि॥\\
२. धर्मराजं तर्पयामि। धर्मराजं तर्पयामि। धर्मराजं तर्पयामि॥\\
३. मृत्युं तर्पयामि। मृत्युं तर्पयामि। मृत्युं तर्पयामि॥\\
४. अन्तकं तर्पयामि। अन्तकं तर्पयामि। अन्तकं तर्पयामि॥\\
५. वैवस्वतं तर्पयामि। वैवस्वतं तर्पयामि। वैवस्वतं तर्पयामि॥\\
६. कालं तर्पयामि। कालं तर्पयामि। कालं तर्पयामि॥\\
७. सर्वभूतक्षयं तर्पयामि। सर्वभूतक्षयं तर्पयामि। सर्वभूतक्षयं तर्पयामि॥\\
८. औदुम्बरं तर्पयामि। औदुम्बरं तर्पयामि। औदुम्बरं तर्पयामि॥\\
९. दध्नं तर्पयामि। दध्नं तर्पयामि। दध्नं तर्पयामि॥\\
१०. नीलं तर्पयामि। नीलं तर्पयामि। नीलं तर्पयामि॥\\
११. परमेष्ठिनं तर्पयामि। परमेष्ठिनं तर्पयामि। परमेष्ठिनं तर्पयामि॥\\
१२. वृकोदरं तर्पयामि। वृकोदरं तर्पयामि। वृकोदरं तर्पयामि॥\\
१३. चित्रं तर्पयामि। चित्रं तर्पयामि। चित्रं तर्पयामि॥\\
१४. चित्रगुप्तं तर्पयामि। चित्रगुप्तं तर्पयामि। चित्रगुप्तं तर्पयामि॥

\closesection

{\sffamily Following this, one should recite the following names as Japa ten times—}

जपः—
\twolineshloka*
{यमो निहन्ता पितृधर्मराजो वैवस्वतो दण्डधरश्च कालः}
{प्रेताधिपो दत्तकृतानुसारी कृतान्तः (एतद् दशकृज्जपन्ति)}

{\sffamily Thereafter, offer Namaskāra—}

नमस्कारः—
\twolineshloka*
{नीलपर्वतसङ्काशो रुद्रकोपसमुद्भवः}
{कालो दण्डधरो देवो वैवस्वत नमोऽस्तु ते}

\closesection
